%=============================================================================!
% 7.5  SYMMETRIES AND CONSERVED QUANTITIES                                    !
%=============================================================================!
\section{Symmetries and conserved quantities}
\label{sec:ha-noether}

Two carts joined by a spring on a straight track (\cref{sec:ne-bodies})
have an energy that depends only on the distance between them, so sliding
both along the track by the same amount changes nothing. Chapter~2 found
that their total momentum is conserved, because the spring pushes them
equally and oppositely. This section shows that the conservation follows
from the symmetry itself, that the energy is unchanged by the shift,
whatever the forces are like, and that the same holds for turning and for
the passing of time.

\subsection*{Moving everything at once}

Let the potential energy of $N$ atoms be unchanged when every atom is
moved by the same displacement $\vb{c}$:
\begin{equation*}
   U(\vb{r}_1 + \vb{c}, \dots, \vb{r}_N + \vb{c})
   = U(\vb{r}_1, \dots, \vb{r}_N)
   \qquad\text{for every } \vb{c} .
\end{equation*}
With $\vb{c} = (c, 0, 0)$ along $x$, we differentiate both sides with
respect to $c$. On the left, each $x_i$ is replaced by $x_i + c$, which
changes at the rate 1 as $c$ grows, so the general chain rule
(\cref{sec:ha-equations}) gives one term for each atom; the right does
not contain $c$. At $c = 0$,
\begin{equation*}
   \sum_{i=1}^{N}\frac{\partial U}{\partial x_i} = 0 ,
   \qquad\text{so}\qquad
   \sum_{i=1}^{N}F_{i,x} = 0 ,
\end{equation*}
and the same for $y$ and $z$: the forces on the atoms add to zero. By
Hamilton's equations $\dot{\vb{p}}_i = \vb{F}_i$, so the total momentum
$\vb{P} = \sum_i\vb{p}_i$ changes at the rate $\sum_i\vb{F}_i = 0$ and is
conserved. Nothing here assumed that the forces act between pairs. The
atoms of \cref{fig:ha-noether}(a) have an energy that depends only on the
angle at the middle one, $U = \frac12 k_\theta(\theta - \theta_0)^2$,
a simple model of the bending of a molecule. The forces,
0.890, 1.430 and \SI{0.685}{\electronvolt\per\angstrom}, are not along the
lines that join the atoms (the force on atom~1 is at right angles to its
bond), yet they add to zero to within $5\times10^{-10}$
\si{\electronvolt\per\angstrom}, the error of the central differences
that computed them. This is the reason promised in \cref{sec:ne-atoms}.

\subsection*{Turning everything at once}

Now let $U$ be unchanged when the whole group is turned by any angle
about any axis through the origin. Turning at the angular velocity
$\vb{n}$, an arrow of length one along the axis, that is at one radian per
unit of time, for a time $\alpha$ turns the group by the angle $\alpha$,
and to first order each atom moves by its velocity times $\alpha$,
$\alpha\,\vb{n}\times\vb{r}_i$ (\cref{sec:ro-rigid}). Differentiating $U$ with respect to
$\alpha$ at $\alpha = 0$ by the general chain rule, one term for each
coordinate of each atom,
\begin{align*}
   0 &= \sum_{i=1}^{N}\grad_iU\cdot(\vb{n}\times\vb{r}_i)
   && \text{since $U$ does not change}\\
   &= \sum_{i=1}^{N}\vb{n}\cdot(\vb{r}_i\times\grad_iU)
   && \text{moving the three vectors round (\cref{sec:ro-tensor})}\\
   &= -\vb{n}\cdot\sum_{i=1}^{N}\vb{r}_i\times\vb{F}_i
   = -\vb{n}\cdot\vb{G}
   && \text{with $\vb{F}_i = -\grad_iU$.}
\end{align*}
This holds for every axis $\vb{n}$, so every component of the total
torque $\vb{G}$ is zero. The angular momentum $\vb{L} = \sum_i\vb{r}_i
\times\vb{p}_i$ changes, by the product rule, at the rate
\begin{equation*}
   \frac{\dd\vb{L}}{\dd t} = \sum_{i=1}^{N}\Bigl(\frac{\vb{p}_i}{m_i}
   \times\vb{p}_i + \vb{r}_i\times\vb{F}_i\Bigr) = \vb{G} = 0 ,
\end{equation*}
since the cross product of an arrow with itself is zero; so $\vb{L}$ is
conserved. For the three atoms of \cref{fig:ha-noether}(a) the torque is
below $2\times10^{-9}$ \si{\electronvolt}, about the origin and about the
point $(3, 2)$~\si{\angstrom}, again the error of the differences. This is the result
promised in \cref{sec:ro-system} for forces that cannot be split into
pairs.

\begin{figure}[htb]
\centering
\includegraphics{ch07_hamilton/noether}
\caption{(a) Three atoms whose energy depends only on the angle at atom~2,
$\frac12 k_\theta(\theta - \theta_0)^2$ with $k_\theta =
\SI{2}{\electronvolt\per\radian\squared}$ and $\theta_0 = \ang{104.5}$,
held at \ang{130} with bonds of 1.0 and \SI{1.3}{\angstrom}. The forces
(arrows, \SI{0.5}{\angstrom} per \si{\electronvolt\per\angstrom}) are not
along the lines joining the atoms (dotted), yet add to zero with no
torque. (b) Two atoms in a box \SI{10}{\angstrom} wide repeated in every
direction, joined by a spring to the nearest copy of each other (ochre);
the forces (\SI{6}{\angstrom} per \si{\electronvolt\per\angstrom}) add to
zero but are not along the line joining the atoms in the box (dotted), so
they exert a torque.}
\label{fig:ha-noether}
\end{figure}

\subsection*{Noether's theorem}

Both results, and the conservation of a momentum whose coordinate is
missing from $\Lag$ (\cref{sec:la-polar}), are cases of one theorem,
proved most directly from the Lagrangian. Suppose that $\Lag$ does not
change, to first order in a small number $\alpha$, when every coordinate
is changed by $\alpha s_k(q)$, with $s_k$ any functions of the
coordinates, and every rate by the matching $\alpha\,\dd s_k/\dd t$, the
rate of $q_k + \alpha s_k$ being $\dot{q}_k + \alpha\,\dd s_k/\dd t$. Then,
by the chain rule,
\begin{align*}
   0 &= \sum_k\Bigl(\frac{\partial\Lag}{\partial q_k}\,s_k
   + \frac{\partial\Lag}{\partial\dot{q}_k}\,\frac{\dd s_k}{\dd t}\Bigr)
   && \text{the change of $\Lag$ per unit of $\alpha$}\\
   &= \sum_k\Bigl(\dot{p}_k\,s_k + p_k\,\frac{\dd s_k}{\dd t}\Bigr)
   && \text{by \cref{eq:la-euler} and $p_k = \partial\Lag/\partial
   \dot{q}_k$}\\
   &= \frac{\dd}{\dd t}\sum_k p_k\,s_k
   && \text{by the product rule,}
\end{align*}
so $\sum_k p_ks_k$ is conserved. Moving every atom along $x$ has $s = 1$
for every $x_i$ and 0 for the other coordinates, and leaves the kinetic
energy unchanged, so $\sum_ip_{i,x}$ is conserved. Turning about
$\vb{n}$ has the change $\vb{n}\times\vb{r}_i$ for atom $i$, and leaves
$\frac12 m_i\lvert\dot{\vb{r}}_i\rvert^2$ unchanged, since the rate
changes by $\alpha\,\vb{n}\times\dot{\vb{r}}_i$, at right angles to
$\dot{\vb{r}}_i$: $\lvert\dot{\vb{r}}_i + \alpha\,\vb{n}\times\dot{\vb{r}}_i
\rvert^2 = \lvert\dot{\vb{r}}_i\rvert^2 + 2\alpha\,\dot{\vb{r}}_i\cdot
(\vb{n}\times\dot{\vb{r}}_i) + \order{\alpha^2}$, and the middle term is
zero (\cref{sec:ro-cross}). So $\sum_i\vb{p}_i\cdot(\vb{n}\times\vb{r}_i)
= \vb{n}\cdot\vb{L}$ is conserved, by the triple product again. A
coordinate missing from $\Lag$ has $s = 1$ for itself and 0 for the
others, and its own momentum is conserved. The passing of time is the
remaining symmetry, though not of this form: \cref{sec:ha-equations}
showed directly that a Hamiltonian that does not contain $t$ is
conserved. Each symmetry of $\Lag$ under a change that can be made as
small as we like, and the passing of time for $H$, comes with a conserved
quantity: this is Noether's theorem, named after Emmy Noether.

\subsection*{A box repeated in every direction}

Most simulations follow the atoms in a box repeated periodically in every
direction (Chapter~10), and each atom feels the nearest copies of the
others. Moving every atom by the same displacement moves every copy too,
and changes nothing, so the forces still add to zero and the total
momentum is conserved. Turning the atoms while the box and the lattice of
its copies stay fixed is not a symmetry. In \cref{fig:ha-noether}(b) two
atoms at $(1, 1)$ and $(9, 2)$~\si{\angstrom} in a box \SI{10}{\angstrom}
wide are joined by a spring of \SI{1}{\electronvolt\per\angstrom\squared}
and natural length \SI{2}{\angstrom} to the nearest copy of each other.
The nearest copy of atom~2, one box width to the left, is at $(9 - 10,
2) = (-1, 2)$~\si{\angstrom}, at the
separation $(-2, 1)$ and the distance $\sqrt5 = \SI{2.2361}{\angstrom}$, so
the spring is stretched by \SI{0.2361}{\angstrom} and pulls atom~1 towards
that copy with $\vb{F}_1 = k\times0.2361\times(-2, 1)/\sqrt5 = (-0.2111,
0.1056)$
\si{\electronvolt\per\angstrom}; atom~2 feels the opposite force, from the
copy of atom~1. The forces add to zero, and moving both atoms by any
displacement leaves the energy at \SI{0.027864}{\electronvolt}. They act along
the line from atom~1 to the copy, at an angle to the line joining the two
atoms in the box, and their torque about the origin is
\begin{equation*}
   G_z = (x_1 - x_2)\,F_{1,y} - (y_1 - y_2)\,F_{1,x}
   = (-8)(0.1056) - (-1)(-0.2111) = \SI{-1.056}{\electronvolt} ,
\end{equation*}
using $\vb{F}_2 = -\vb{F}_1$ in $G_z = \sum_i(x_iF_{i,y} - y_iF_{i,x})$.
Turning both atoms about the origin changes the energy at the rate
$\SI{1.056}{\electronvolt}$ per radian, so the angular momentum of the
atoms in the box is not conserved, and \cref{sec:ro-remove} removed only
the drift of a periodic system.

A learned model of the forces keeps these conservation laws when its
forces are the slopes of an energy with the same symmetries. A learned
model of the motion predicts the next state without Hamilton's equations,
so the theorem does not cover it even when it respects the symmetries,
which is why TrajCast corrects its predicted momenta, as the planned
Chapter~28 will discuss.

\begin{keyidea}
Every symmetry of the Lagrangian under a change that can be made as small
as we like gives a conserved quantity (Noether's theorem), and so does
the passing of time: unchanged by
moving everything, the total momentum; by turning everything, the angular
momentum; by the passing of time, the energy. No assumption about pairs of
forces is needed.
\end{keyidea}

\notebook{07}{forces of a bending molecule and of atoms in a periodic box}

\begin{selfcheck}
A fixed flat wall at $x = 0$ that repels atoms is added to a periodic
box. Which components of the total momentum are still conserved?
\end{selfcheck}
